\section{Related work}

\subsection{Significance of Tree Species Classification and Conventional Approaches}
Tree species classification in urban environments serves as a fundamental technology for landscape management and conservation planning \cite{ref4}. In recent years, automatic identification using smartphone applications has become increasingly common. Several mobile applications have been developed, including LeafSnap \cite{ref5} , Pl@ntNet \cite{ref6} , Folia \cite{ref7}. These tools are widely used by citizens and in educational contexts, providing an accessible platform for understanding urban biodiversity, even for non-specialists. However, most of these applications rely primarily on close-up images, and research focusing on mid- or long-range observations remains limited. Furthermore, many existing methods require manual annotation for supervised classification, which leads to high labeling costs in practical use. 

\subsection{Commonly Used Features and Limitations of Existing Datasets}
Typical morphological features used for tree classification include leaves, bark (trunks), leaf shapes, flowers, and fruits. Among these, leaves are the most commonly used because they are easy to observe. Large-scale datasets such as Pl@ntLeaves \cite{ref8} and LeafSnap \cite{ref5} have contributed to the prevalence of leaf-based classification studies. However, leaves, flowers, and fruits are highly seasonal and cannot be used during defoliation or non-flowering periods. In contrast, bark can be observed year-round and is thus a promising feature independent of seasonal variations. Public datasets such as BarkNet 1.0 \cite{ref9} and Bark-101 \cite{ref10} have facilitated progress in bark-based classification research. Nevertheless, most of these datasets consist of close-up images captured under conditions with minimal background interference. When relying solely on texture features, differences between bark textures of various species can be subtle, and classification becomes even more challenging as trees age. Moreover, when trunks are photographed from mid- to long-range distances, variations in shooting distance and image resolution significantly affect texture representation, leading to accuracy degradation—commonly referred to as the texture degradation problem. Consequently, the applicability of bark-based classification remains limited under realistic conditions.

\subsection{Urban Tree Datasets and Future Directions}
Most existing tree datasets have been independently collected by individual researchers, resulting in limited interoperability and comparability between datasets \cite{ref4}. In contrast, the UrbanStreetTree dataset \cite{ref11}, developed for urban street trees in China, covers multiple tree organs—including trunks, leaves, whole trees, branches, flowers, and fruits—and contains 41,467 images. It can be used for both classification and segmentation tasks. A distinctive feature of this dataset is its detailed annotation of branch structures, which enables species identification even under challenging conditions such as winter leaf-off periods or when trunks are painted.

\subsection{Positioning of This Study}
To address the above challenges, this study aims to reduce background dependency by combining the ViT \cite{ref1} with ABMIL \cite{ref2}, encouraging the model to focus on trunk regions. The MIL-based patch division allows the model to retain high-resolution local information while avoiding information loss due to extensive resizing, thereby facilitating attention to fine-grained features even in long-distance images. Furthermore, by leveraging the self-attention mechanism of the Transformer architecture, the proposed framework can simultaneously capture both local bark texture and global tree structure, leading to more robust and context-aware tree classification.

